\documentclass[10pt,a4paper]{ctexart}
\usepackage[margin=2cm]{geometry}
\usepackage{graphicx}
\usepackage{float}      % 支持 [H] 浮动体选项
\usepackage{hyperref}
\usepackage{amsmath}    % 提供方程环境
\usepackage{physics}    % 简化向量、微分算子语法
\newcommand{\myspace}[1]{\par\vspace{#1\baselineskip}} %自定义空一行
\usepackage{fancyhdr}
\pagestyle{fancy}
\fancyhf{}  % 清除所有页眉页脚
\fancyfoot[C]{\thepage}  % 在页脚中央设置页码
\renewcommand{\headrulewidth}{0pt}  % 去掉页眉横线
\renewcommand{\footrulewidth}{0pt}  % 去掉页脚横线

\usepackage{hyperref}
\hypersetup{
    colorlinks=true,
    linkcolor=blue,
    filecolor=magenta,      
    urlcolor=cyan,
    pdftitle={Overleaf Example},
    pdfpagemode=FullScreen,
    }

\usepackage{titlesec}

% 重新定义section格式
\titleformat{\section}
  {\normalfont\Large\bfseries}
  {\chinese{section}、}
  {0pt}
  {}

% 给出常用字号的自定义指令
\newcommand{\chuhao}{\fontsize{42pt}{48pt}\selectfont}
\newcommand{\xiaochu}{\fontsize{36pt}{42pt}\selectfont}
\newcommand{\xiaoer}{\fontsize{18pt}{21.6pt}\selectfont}
\newcommand{\sanhao}{\fontsize{16pt}{20pt}\selectfont}
\newcommand{\xiaosan}{\fontsize{15pt}{18pt}\selectfont}
\newcommand{\sihao}{\fontsize{14pt}{18pt}\selectfont}
\newcommand{\xiaosi}{\fontsize{12pt}{16pt}\selectfont}
\newcommand{\wuhao}{\fontsize{10.5pt}{12.6pt}\selectfont}

% 定义中文数字转换
\titleformat{\section}
  {\normalfont\Large\bfseries}
  {\zhnum{section}、}  % 使用\zhnum而不是\chinese
  {0pt}
  {}

% 标题设置
\title{\xiaochu\textbf{物理实验预习报告}}
\author{}
\date{}

\begin{document}

\vspace*{36pt} % 顶部空白

% 居中填写区域
\begin{center}   
    \xiaochu\textbf{物理实验预习报告}
    % 预留空白区域
    \vspace{13pt}
    \vspace{13pt}
    \vspace{13pt}
    \vspace{13pt}
    
    \sanhao\textbf{实验名称：\underline{\makebox[10cm]{ 万用表的设计}}} \\[10pt]
    \sanhao\textbf{指导教师：\underline{\makebox[10cm]{ 唐浩原 }}} \\[30pt]
    
    % 预留空白区域
    \myspace{6}
    
   \sihao\textbf{班级：\underline{\makebox[6cm]{}}} \\[10pt]
    \sihao\textbf{姓名：\underline{\makebox[6cm]{}}} \\[10pt]
    \sihao\textbf{学号：\underline{\makebox[6cm]{}}} \\[30pt]
    
   \sihao\textbf{实验日期: 
   \underline{\makebox[0.8cm]{2026}}  年 
   \underline{\makebox[0.4cm]{1}}月
   \underline{\makebox[0.4cm]{5}}日  
   星期\underline{\makebox[0.6cm]{三}}上午}
    
    \vspace{18pt}
    \sihao 浙江大学物理实验教学中心
    
\end{center}

\newpage

    \section{实验综述}
    
    \xiaosi （自述实验现象、实验原理和实验方法，不超过300字，5分）

本实验主要包括电表内阻测量及电表改装两部分内容。首先采用替代法和中值法精确测定微安表的内阻，这是后续电表改装的基磁基础。随后进行电表量程扩展：通过在表头两端并联适当的分流电阻，可将微安表改装为毫安表或安培表；而在表头上串联合适的分压电阻，则可将其改装为伏特表。最后，利用改装的电流表配合电源和调零电阻构成欧姆表，用以测量未知电阻。当欧姆表示数为$I_x$时，设其等效内阻为$R_0$，电源电动势为$E$，根据闭合电路欧姆定律可得：
$$ E = I_x(R_0 + R_x) $$
由此解得待测电阻值为$R_x = \dfrac{E}{I_x} - R_0$。整个实验过程需注重注重电路连接的正确性及数据的准确记录与分析。
\myspace{2}
    
    \section{实验重点}
    
    \xiaosi（简述本实验的学习重点，不超过100字，3分）

掌握电表内阻测量的替代法与中值法原理，理解电表改装的电路设计方法。学会计算分流电阻和分压电阻的阻值，完成电流表和电压表的改装。理解欧姆表的测量原理，能够分析改装后电表的刻度特性及测量误差来源。

    \myspace{2}
    
    \section{实验难点}

    \xiaosi（简述本实验的实现难点，不超过100字，2分）

    电表内阻的精确测量是首要难点，需注意测量过程中对表头的保护。分流电阻和分压电阻的计算与选取要求精准，否则会引入较大系统误差。欧姆表中值电阻的确定及刻度非线性的补偿也是实际操作中的技术难点。
\myspace{2}


\newpage

\sihao\textbf{注意事项：}
\xiaosi
\begin{enumerate}
    \item 用PDF格式上传“预习报告”，文件名：学生姓名+学号+实验名称+周次。
    \item “预习报告”必须递交在“学在浙大”的本课程的对应实验项目的“作业”模块内。
    \item “预习报告”还须拷贝到“实验报告”中（便以教师批改）。
    \item 大学物理实验”课程选做实验使用本“预习报告”；必做实验无须写预习报告，前往“学在浙大”完成预习测试即可。
\end{enumerate}

\vspace{1cm}
\xiaosi\centering \textbf{浙江大学物理实验教学中心制}

\end{document}
